% Nouveau EGA V (cote 25): archive cote 25, site page 179 (PDF page 180).
% typescript. Draft compared: none.
% AI-transcribed by Claude Opus 5.5 (claude-opus-5-5, adaptive thinking, effort high), Message Batches API, batch ega25b (msgbatch_01A7n5YuXubUPNK1zsepi8RG), 2026-11-13; not yet checked by eye.
\scanpage{25/179}
\mspage{\hand{6}}

\[ \underline{Q}(\underline{E}) =\xout{W}(\underline{Sym}^2(\underline{E})) \]
et définissons de même $\underline{Q}(\underline{E}_o)$; par l'opération de ``transport de
structure'' on aura donc un morphisme $u \rightsquigarrow \underline{Q}(u)$ :
\[ \underline{Isom}(\underline{E}_o,\underline{E}) \longrightarrow \underline{Isom}(\underline{Q}(\underline{E}_o),\underline{Q}(\underline{E})), \]
et, utilisant la section $q_o$ de $\underline{Q}(\underline{E}_o)$ correspondant à $Q_o$ , on trouve un
morphisme $u \rightsquigarrow \underline{Q}(u)(q_o)$:
\[ (*) \qquad \underline{Isom}(\underline{E}_o,\underline{E}) \longrightarrow \underline{Q}(\underline{E}) \quad . \]
D'autre part, la forme Q \xout{\unsure{sur E}} correspond à une section q de $\underline{Q}(\underline{E})$ sur S,
et P n'est autre que l'image inverse de q(S) par (*), comme il résulte
trivialement des définitions, (ce qui, d'ailleurs, établit en passant la
représentabilité annoncée de P \xout{z}comme préschéma affine et de \xout{type}\inserted{prés} finie
sur S). Introduisons maintenant l'ouvert \xout{$\underline{Q}$x}
\[ \underline{Q}(\underline{E})^* = \underline{Q}(\underline{E})_{\tilde{d}} \]
de $\underline{Q}(\underline{E})$ , correspondant aux formes quadratiques de discriminant corrigé
inversible (représentan\xout{x}t le foncteur $S' \rightsquigarrow$ ensemble des \xout{formes}
sections de $\underline{Sym}^2(\underline{E}_{S'})$ à discriminant corrigé inversible). En fait,
par transport de structure on a $\underline{Isom}(\underline{E}_o,\underline{E}) \rightarrow \underline{Isom}(\underline{Q}(\underline{E}_o)^*,\underline{Q}(\underline{E})^*)$, et
comme $q_o$ est une section de $\underline{Q}(\underline{E}_o)^*$ (puisque le discriminant \unsure{corrigé|modrigé} de
la forme type $Q_o$ vaut 1), le morphisme (*) se factorise en fait en un
morphisme
\[ (*\,*) \qquad \underline{Isom}(\underline{E}_o,\underline{E}) \rightarrow \underline{Q}(\underline{E})^* , \]
qui est évidemment de présentation finie, puisque les deux membres sont de
présentation finie sur S. Il suffit maintenant de prouver la
\xout{Lemme} \uline{Proposition} 16.11. Avec les notations précédentes (à rappeler),
le morphisme (* *) est \uline{fidèlement plat} .

Il s'ensui\deleted{t}\inserted{vra} que si Q satisfait à (ii), i.e. \inserted{si} q est en fait une sectio[n]
de $\underline{Q}(\underline{E})^*$\hand{\unsure{$\hat{}$}} , alors P (déduit de (**) par \inserted{\unsure{le}} changement de base q) est encore

\mspage{\hand{179}}
